\section{\solution~Design}
\label{sec:design}

Fig.~\ref{fig:arch} depicts the components of \solution. It is composed of a runtime binding that ties each tool call to its execution, a kernel collector, a fusion resolver, a live accounting daemon, and a per-execution egress policy that uses the same~identity.

\begin{figure}[!h]
\centering
\begin{tikzpicture}[font=\scriptsize, x=1cm, y=1cm,
  comp/.style={draw, fill=white, align=center, inner sep=2pt, minimum height=5.5mm},
  env/.style={draw, fill=gray!30, align=center, inner sep=2pt},
  arr/.style={-{Stealth[length=1.6mm]}, semithick},
  lab/.style={font=\tiny, inner sep=1pt}]
\draw[fill=gray!8] (0,-0.25) rectangle (8.6,5.0);
\draw (0,2.4) -- (8.6,2.4);
\node[rotate=90] at (0.25,3.7) {User space};
\node[rotate=90] at (0.25,1.0) {Kernel};
\node[env, minimum width=2.4cm, minimum height=1.45cm, anchor=north] (rt) at (2.05,4.85) {};
\node[anchor=north, align=center] at (rt.north) {Agent runtime\\(MCP server)};
\node[comp, minimum width=1.3cm, minimum height=3.5mm, inner sep=1pt] at (2.17,3.73) {};
\node[comp, minimum width=1.3cm, minimum height=3.5mm, inner sep=1pt] at (2.11,3.67) {};
\node[comp, minimum width=1.3cm, minimum height=3.5mm, inner sep=1pt] (tc) at (2.05,3.61) {Tool calls};
\node[comp, fill=gray!55, font=\scriptsize\bfseries, minimum width=2.8cm, minimum height=1.0cm] (la) at (6.35,4.25) {Live Accounting};
\draw[rounded corners=6pt, fill=gray!55] (0.75,0.55) rectangle (8.3,3.2);
\node[font=\scriptsize\bfseries, anchor=south west] at (0.85,0.6) {\solution};
\node[comp, minimum width=2.0cm] (tsb) at (2.05,2.85) {Task-Scoped Binding};
\node[comp, minimum width=2.6cm] (fr) at (6.35,2.85) {Fusion Resolver};
\node[comp, minimum width=1.75cm] (cb) at (2.05,1.6) {Context Binding};
\node[comp, minimum width=1.5cm] (col) at (4.25,1.6) {Collector};
\node[comp, minimum width=1.95cm, align=center] (pol) at (6.6,1.6) {Per-Execution\\Egress Policy};
\node[env, minimum width=3.6cm, minimum height=4.5mm] (tcp) at (5.55,0.12) {TCP/IP stack (\texttt{connect})};
\draw[arr] (tc) -- (tsb);
\draw[arr] (tsb) -- node[lab, right]{identity} (cb);
\draw[arr] (cb) -- (col);
\draw[arr] (cb.south) -- ++(0,-0.38) -| ([xshift=-4mm]pol.south);
\draw[arr] (col.north) |- (fr.west);
\draw[arr] (fr) -- (la);
\draw[arr] (tcp.north -| col) -- (col.south);
\draw[arr] ([xshift=4mm]pol.south) -- node[lab, right]{allow / deny} ([xshift=4mm]pol.south |- tcp.north);
\begin{scope}[shift={(0.3,-0.65)}, every node/.style={font=\tiny, anchor=west}]
\draw[fill=gray!55] (0,0) rectangle +(0.25,0.18); \node at (0.3,0.09) {\solution{} components};
\draw[fill=gray!30] (2.6,0) rectangle +(0.25,0.18); \node at (2.9,0.09) {agent environment and host stack};
\draw[fill=gray!8] (5.9,0) rectangle +(0.25,0.18); \node at (6.2,0.09) {host server};
\end{scope}
\end{tikzpicture}
\caption{\solution{} architecture.}
\label{fig:arch}
\end{figure}

\subsection{Collector}

eBPF programs record, for every thread, \texttt{sched\_process\_fork}, \texttt{exec}, and \texttt{exit}, from which lineage is reconstructed. They also hook \texttt{tcp\_connect} in process context, capturing PID, thread ID, cgroup v2 ID, current context, and socket. The handshake outcome is taken from \texttt{inet\_sock\_set\_state} and joined by socket. Socket state changes may run in softirq context, so the PID observed there is never used for attribution.

\subsection{Context Binding}

The runtime hands the execution identity to the kernel through markers that the collector records in a kernel map keyed by thread ID; on fork or thread creation, the child inherits the parent's entry, so a tool that runs as a subprocess carries the identity without any change to the tool. In the prototype, a marker is an \texttt{openat} call on a reserved, nonexistent path, captured by a tracepoint; it needs no library and costs one failed system call. Emitting markers only at the start and end of each tool call (\emph{manual binding}) is correct for blocking, threaded, and subprocess tools, but not for asynchronous ones: coroutines of different executions interleave on the event-loop thread and overwrite its entry (Section~\ref{sec:results}).

\subsection{Task-Scoped Binding}
\texttt{fact\_rt} makes the binding follow the task instead of the thread. The execution identity is stored in a context variable, which the asynchronous runtime copies into every task and into worker threads. In Python, the library wraps \texttt{socket.connect}/\texttt{connect\_ex} and process creation (\texttt{subprocess.Popen} and \texttt{os.fork}) and, immediately before the system call, sets the calling thread's kernel entry to the identity of the current task, or clears it if the task has none. Because no other coroutine can run between the marker and the system call, each connection and each child process carries the identity of the task that created it, regardless of interleaving. The entry is not cleared after the call; a connection that bypasses these wrappers (\eg from native code) inherits the identity of the last task that connected on that thread, a limitation we discuss in Section~\ref{sec:discussion}. For MCP servers, a wrapper around the \texttt{tools/call} handler takes the identity from the request's \texttt{\_meta} field, which the protocol reserves for metadata, and binds it for the duration of the call, so tool code is unmodified.

\subsection{Fusion Resolver}
For each connection the resolver receives only $(t, \mathit{pid}, \mathit{tid}, \mathit{cgroup}, \mathit{ctx})$ and computes three candidate sets: $L$ from the nearest registered lineage root, $G$ from the cgroup registry, and $c$ from the context binding. Let $S$ be the intersection of the non-empty structural sets, or their union if they are disjoint (structural conflict). The output is:
\begin{itemize}
  \item $\{c\}$ if a context is present and $S=\emptyset$ or $c \in S$;
  \item $S \cup \{c\}$ (ambiguous) if a context is present and $c \notin S$, \ie the context contradicts the process structure;
  \item $S$ otherwise, which may be a single origin, a candidate set, or empty.
\end{itemize}
The resolver never selects among candidates that the available information cannot separate. When the structure is informative (distinct trees), a mislabeled context therefore becomes ambiguity rather than misattribution. When the structure is uninformative (a shared worker), correctness rests entirely on the binding.

\subsection{Live Accounting}
\texttt{factd} consumes the collector stream, reorders events from different CPUs within a short window, applies the same fusion function as the offline evaluation, and periodically exports one flow-like record per execution, in the format of IP Flow Information Export (IPFIX)~\cite{rfc7011} (\ie connections, established, failed, distinct destinations), plus the ambiguous and unattributed counts with their candidate sets. Ambiguous or unattributed connections are never credited to an execution.

\subsection{Per-Execution Egress Policy}

A \texttt{cgroup/connect4} eBPF program, loaded with libbpf, keeps its own copy of the per-thread identity map (fed by the same markers) and decides each Internet Protocol version 4 (IPv4) \texttt{connect} by looking up the pair (execution, destination) in an allow-list map; connections without identity receive a configurable default action (\textit{deny} in our experiments), instead of inheriting a guess. Every decision is logged to a ring buffer for auditing. Thus, policy and accounting share one notion of identity, and errors in the binding become measurable policy errors: a false allow is a policy bypass and a false \textit{deny} is a loss of availability. The policy also applies the resolver's lineage check in the kernel. The launcher registers the root of every process tree that belongs to a single execution, and the root is inherited on fork and thread creation. When a thread's label contradicts the owner of its tree, the default mode denies the connection and logs a conflict; an alternative mode decides by the tree owner. A tree shared by several executions (a shared worker) is not registered, so there the decision follows the label.
